\documentclass[11pt]{article}
\usepackage{manual_docs_style}
\externaldocument{build/terminology_shared_utilities}[terminology_shared_utilities.pdf]
\externaldocument{build/appendix_environment}[appendix_environment.pdf]

\begin{document}

\newcommand{\environmentBuildMetadataStageLink}{%
  \ifdefined\tracebenchDocsCombinedManualMode
    \hyperref[sec:simulink-build-metadata-stage]{build metadata stage}%
  \else
    \href{simulink_generated_metadata_schema.pdf}{build metadata stage}%
  \fi
}

\docTitle{Environments}
\phantomsection\label{docstart:environment_profiles}

\section*{Overview}

This document describes the manually created files under \code{models/simulink/<simulator_id>/} that define a \name{} \termSimulator.
These files, together with  \code{simulink_model_original.mdl}, define a simulator environment and are necessary to generate candidate runs, simulate them, add noise, and render prompts. 
\phantomsection\label{sec:allowed-tracebench-models}
Additionally, \termSimulator needs to be added to \hyperref[sec:allowed-tracebench-models]{\code{shared/benchmark_utils.py::ALLOWED_TRACEBENCH_MODELS}} and the \environmentBuildMetadataStageLink{} needs to be run.


The files are located in \code{models/simulink/<simulator_id>/}.
\newcommand{\environmentTreeSimulinkModelOriginal}{\hyperref[sec:environment-simulink-model-original]{simulink\_model\_original.mdl}}
\newcommand{\environmentTreeDescriptionLevels}{\hyperref[sec:environment-description-levels]{description\_levels.json}}
\newcommand{\environmentTreeExperimentConfig}{\hyperref[sec:environment-experiment-config]{experiment\_config.json}}
\newcommand{\environmentTreeNoiseAdder}{\hyperref[sec:environment-noise-adder]{noise\_adder.py}}
\newcommand{\environmentTreeGenerationPolicy}{\hyperref[sec:generation-policy-schema]{generation\_policies/\textless policy\_id\textgreater.json}}
\newcommand{\environmentTreeDetectabilitySpecific}{\hyperref[sec:environment-detectability-specific]{detectability\_specific\_environment.py}}
\begin{LinkedDocCode}
models/simulink/<simulator_id>/
|-- \environmentTreeSimulinkModelOriginal
|-- \environmentTreeDescriptionLevels
|-- \environmentTreeExperimentConfig
|-- \environmentTreeNoiseAdder
|-- \environmentTreeGenerationPolicy
|-- \environmentTreeDetectabilitySpecific\quad
\end{LinkedDocCode}

In the following each of these files is described.
\phantomsection\label{sec:generation-policy-schema}
\subsection*{Generation Policy Schema}


\paragraph{Policy Generation}
Note that the original public test release on Hugging Face did not use this generation-policy-based sampling pipeline. 
After introducing generation policies, we backfilled a 10k-family candidate pool for each simulator from the original dataset release, using the same sampling distribution used to construct the test set.

Since this way of sampling was introduced after the first public released of the dataset, we have a 10k version for each of the \termSimulator in the first version in distribution with the test set. 


The runs are sampled a generation policy.

A generation policy defines:

\begin{itemize}
\item how many candidate families to generate;
\item how to sample baseline parameters and initial states;
\item how many interventions to sample, and how to sample intervention times and post-intervention values.
\end{itemize}

Candidate generation resolves the policy into a concrete plan before simulation.

\subsubsection*{Top-Level Fields}

Required top-level fields in a generation policy. All fields listed in this table are required.

\begin{center}
\begin{tabular}{p{0.28\linewidth}p{0.60\linewidth}}
\hline
\textbf{Field} & \textbf{Meaning} \\
\hline
\code{policy_id} & Non-empty identifier for the policy. Should match the filename stem. \\
\code{model} & \termSimulator identifier under \code{models/simulink/<simulator_id>}. \\
\code{seed} & Integer seed used to make candidate generation reproducible. \\
\code{target_candidate_families} & Number of baseline-centered candidate families to sample. \\
\code{baseline_sampler} & Defines how baseline parameters and initial states are sampled. \\
\code{intervention_sampler} & Defines intervention parameters, directions, times, and values. \\
\hline
\end{tabular}
\end{center}

\subsubsection*{Minimal Shape}

\begin{DocCode}
{
  "policy_id": "production_v1",
  "model": "BallDrop",
  "seed": 456,
  "target_candidate_families": 1000000,
  "baseline_sampler": {
    "method": "latin_hypercube"
  },
  "intervention_sampler": {
    "parameters": "all",
    "directions": ["increase", "decrease"],
    "children_per_parameter_direction": 3,
    "intervention_time": {
      "distribution": "uniform",
      "min_fraction_of_horizon": 0.2,
      "max_fraction_of_horizon": 0.8
    },
    "value_sampler": {
      "method": "conditional_uniform",
      "distance_space": "sampling",
      "min_abs_distance_multiplier": 1.0,
      "min_srd_distance_multiplier": 1.0
    }
  }
}
\end{DocCode}

\subsubsection*{\texorpdfstring{\titlecode{baseline_sampler}}{baseline\_sampler}}

\code{baseline_sampler} defines how baseline parameters and initial-state variables are sampled.
It samples over the variables listed in:

\begin{DocCode}
experiment_config.json::exposed_variables.parameters
experiment_config.json::exposed_variables.initial_state
\end{DocCode}

The sampler must respect all \code{allowed_intervals} defined in \code{experiment_config.json}.

\paragraph{Schema}

\begin{DocCode}
"baseline_sampler": {
  "method": "latin_hypercube",
  "variables": {
    "mass": {
      "distribution": "loguniform"
    },
    "gravity": {
      "distribution": "uniform"
    },
    "initial_height": {
      "distribution": "uniform"
    }
  },
  "baseline_min_distance_scale": 0.2
}
\end{DocCode}

\paragraph{Fields}

\begin{itemize}
\item \code{method}: required string. Allowed values are \code{independent}, \code{latin_hypercube}, \code{sobol}, and \code{grid}.
\item \code{variables}: optional object keyed by exposed variable name. It overrides the default distribution for selected variables.
\item \code{baseline_min_distance_scale}: optional non-negative number. If present, generated baselines should be separated from one another according to the policy's diversity constraint.
\end{itemize}

Allowed variable-level distributions are \code{uniform,loguniform,fixed}.
For \code{uniform} and \code{loguniform}, an optional
\code{interval = [lower, upper]} may condition sampling on a strict sub-interval
of the variable's environment-level \code{allowed_intervals}.  The conditioned
interval must be finite, increasing, and entirely contained in the
environment-level interval; it does not change the environment contract.

\subsubsection*{\texorpdfstring{\titlecode{intervention_sampler}}{intervention\_sampler}}

\code{intervention_sampler} defines how intervention children are generated from each baseline family.

\paragraph{Schema}

\begin{DocCode}
"intervention_sampler": {
  "parameters": "all",
  "directions": ["increase", "decrease"],
  "children_per_parameter_direction": 3,
  "intervention_time": {
    "distribution": "uniform",
    "min_fraction_of_horizon": 0.2,
    "max_fraction_of_horizon": 0.8
  },
  "value_sampler": {
    "method": "conditional_uniform",
    "distance_space": "sampling",
    "min_abs_distance_multiplier": 1.0,
    "min_srd_distance_multiplier": 1.0
  }
}
\end{DocCode}

\paragraph{Fields}

\begin{itemize}
\item \code{parameters}: \code{"all"} or a non-empty list of names from \code{experiment_config.json::exposed_variables.parameters}; \code{"all"} selects all listed parameters.
\item \code{directions}: non-empty list containing \code{"increase"}, \code{"decrease"}, or both.
\item \code{children_per_parameter_direction}: positive integer.
\item \code{intervention_time}: object defining how intervention times are sampled.
\item \code{value_sampler}: object defining how post-intervention values are sampled.
\end{itemize}

Only variables listed in \code{experiment_config.json::exposed_variables.parameters} may be intervened on.
Variables listed under \code{initial_state} are baseline-only variables and must not change after simulation start.

\paragraph{\texorpdfstring{\titlecode{intervention_time}}{intervention\_time}}

\code{intervention_time} defines how the step-change time is sampled.

\begin{DocCode}
"intervention_time": {
  "distribution": "uniform",
  "min_fraction_of_horizon": 0.2,
  "max_fraction_of_horizon": 0.8
}
\end{DocCode}

Allowed fields:

\begin{itemize}
\item \code{distribution}: required string. Allowed values are \code{uniform}, \code{fixed}, and \code{choice}.
\item \code{min_fraction_of_horizon}: required for \code{uniform}. Float in \code{[0,1]}.
\item \code{max_fraction_of_horizon}: required for \code{uniform}. Float in \code{[0,1]}.
\item \code{value}: required when \code{distribution} is \code{fixed}.
\item \code{values}: required when \code{distribution} is \code{choice}.
\end{itemize}

All intervention times must satisfy:

\begin{DocCode}
0 < intervention_time < experiment_config.json::end_time_input_s
\end{DocCode}

If \code{min_fraction_of_horizon} and \code{max_fraction_of_horizon} are used, the actual interval is computed from \code{experiment_config.json::end_time_input_s}.

\paragraph{\texorpdfstring{\titlecode{value_sampler}}{value\_sampler}}

\code{value_sampler} defines how the post-intervention value is sampled for each changed parameter.

\begin{DocCode}
"value_sampler": {
  "method": "conditional_uniform",
  "distance_space": "sampling",
  "min_abs_distance_multiplier": 1.0,
  "min_srd_distance_multiplier": 1.0,
  "max_srd_distance_multiplier": 1.2
}
\end{DocCode}

Allowed \code{method} values:

\begin{itemize}
\item \code{conditional_uniform}: uniformly sample valid values on the requested side of the baseline value after applying distance constraints.
\item \code{fixed_delta}: add or subtract a fixed amount. Requires \code{delta}.
\item \code{fixed_srd}: sample values targeting a fixed SRD distance. Requires \code{target_srd}.
\end{itemize}

\code{distance_space} is optional and accepts \code{linear} or \code{sampling}; the default is \code{linear} for backward compatibility.
For a \code{loguniform} variable, \code{sampling} applies the SRD distance constraint in logarithmic sampling space, while the minimum absolute-distance constraint remains in the variable's original units.
For other distributions the two modes are equivalent.
\code{max_srd_distance_multiplier} is optional.  When present, it places an
upper bound on the sampled distance in the selected distance space and must be
greater than or equal to \code{min_srd_distance_multiplier}.  It is useful for
targeted yield pilots that must avoid needlessly extreme changes while retaining
the environment-level minimum distance.

The generated post-intervention value must satisfy all of the following:

\begin{itemize}
\item it lies inside \code{experiment_config.json::exposed_variables.parameters::<parameter>::allowed_intervals};
\item it changes exactly one root-cause parameter;
\item its absolute distance from the baseline value is at least the environment-level \termMinAbsDist{} multiplied by \code{min_abs_distance_multiplier};
\item its SRD distance from the baseline value is at least the environment-level \termMinSRDDistance{} multiplied by \code{min_srd_distance_multiplier};
\item when \code{max_srd_distance_multiplier} is present, its distance in the
selected distance space is at most the environment-level
\termMinSRDDistance{} multiplied by that value.
\end{itemize}

If no valid value exists for a requested parameter and direction, the generator must not create a runnable child node.
The missing direction must instead be recorded in \code{generation_report.json}.

\phantomsection\label{sec:environment-description-levels}
\subsection*{\texorpdfstring{\titlecode{description_levels.json}}{description\_levels.json}}

\code{description_levels.json} contains \termSimulator-specific text used by the prompt renderer.
It has:
\begin{itemize}
\item \code{textual_description}: object keyed by \termDescLevel.
\begin{itemize}
\item \code{high}: detailed natural-language description, including relevant equations, observed channels, and root-cause parameters.
\end{itemize}
\item \code{observed_signals_description}: signal description used in prompts.
\item \code{internal_naming_to_agent_facing_signal}: object mapping internal signal names to agent-facing signal descriptions.
The non-\code{time} keys must match \code{experiment_config.json::observable_signals::observable_signals} exactly and in the same order.
\begin{DocCode}
{
  "signal_internal_name_0": "agent-facing description 0",
  "signal_internal_name_1": "agent-facing description 1",
  "time": "time"
}
\end{DocCode}
\item \code{internal_naming_to_agent_facing_parameter}: object mapping internal intervenable parameter names from \code{experiment_config.json::exposed_variables.parameters} to agent-facing labels.
\end{itemize}


\phantomsection\label{sec:environment-experiment-config}
\subsection*{\texorpdfstring{\titlecode{experiment_config.json}}{experiment\_config.json}}

\code{experiment_config.json} defines the environment contract.
It should contain:

\begingroup
\footnotesize
\renewcommand{\arraystretch}{1.12}
\setlength{\tabcolsep}{4pt}
\begin{longtable}{|P{0.27\textwidth}|P{0.18\textwidth}|P{0.45\textwidth}|}
\caption{Experiment config top-level fields.}\\
\hline
\textbf{Name} & \textbf{Type} & \textbf{Description} \\
\hline
\endfirsthead
\hline
\textbf{Name} & \textbf{Type} & \textbf{Description} \\
\hline
\endhead
\code{paper_facing_name} & string & Environment name used in the paper and generated tables. \\
\hline
\code{exposed_variables} & object & Variables externally exposed to \name{}. \\
\hline
\code{observable_signals} & object & Non-time channels observed by agents and metrics. \\
\hline
\code{sampling_rate_hz} & positive number & Target sampling rate for materialized trajectories. \\
\hline
\code{end_time_input_s} & positive number & Simulation horizon. \\
\hline
\code{detectability} & object & Thresholds used by Stage: Cheap Filter Metrics. \\
\hline
\end{longtable}
\endgroup
\agentProposedEdit
  {agent-add-3e3ae322-e7bf-5762-b638-e99a444010eb}
  {Add this agent-authored block for review.}
  {}
  {The optional top-level boolean \code{allow_family_baseline_reuse} determines whether multi-worker simulation may use Fast Restart family-baseline reuse and defaults to \code{true} when omitted.
When \code{allow_family_baseline_reuse} is \code{false}, baselines and interventions use the individual simulation path even when multiple MATLAB workers are available.}

\subsubsection*{\code{exposed_variables}}

\code{exposed_variables} has two groups:

\begin{itemize}
\item \code{parameters}: root-cause parameters that may be changed by an intervention.
\item \code{initial_state}: baseline-only state variables that are sampled at the start of a run and must not change after simulation start.
\end{itemize}

Each exposed variable should define:

\begin{itemize}
\item \code{allowed_intervals}: legal range \code{[min, max]}.
\item \termMinSRDDistance
\item \termMinAbsDist
\end{itemize}

\subsubsection*{\code{observable_signals}}

\code{observable_signals} defines the channels that agents and downstream metrics see.

\begin{itemize}
\item \code{observable_signals}: ordered list of non-time signal names.
\item \code{signal_type}: object keyed by signal name. Each value has \code{type}, which is either \code{continuous} or \code{impulse_like}.
\end{itemize}

\subsubsection*{\code{detectability}}

\code{detectability} contains the thresholds used by Stage: Cheap Filter Metrics to decide whether an intervention is visible in the observed signals.
It is divided into continuous-signal rules, impulse-like-signal rules, and per-signal RMS thresholds.

\begingroup
\footnotesize
\renewcommand{\arraystretch}{1.12}
\setlength{\tabcolsep}{4pt}
\begin{longtable}{|P{0.20\textwidth}|P{0.28\textwidth}|P{0.42\textwidth}|}
\caption{Detectability fields.}\\
\hline
\textbf{Block} & \textbf{Field} & \textbf{Meaning} \\
\hline
\endfirsthead
\hline
\textbf{Block} & \textbf{Field} & \textbf{Meaning} \\
\hline
\endhead
\code{continuous} & \code{min_srd_distance} & Minimum \termSRD distance used for threshold-crossing detection on continuous signals. \\
\hline
\code{continuous} & \code{epsilon_SRD} & Epsilon term used in the \termSRD denominator. \\
\hline
\code{continuous} & \code{minimum_consecutive_below_SRD} & Number of consecutive samples that must exceed the \termSRD threshold before a difference is recorded. \\
\hline
\code{impulse_like} & \code{min_srd_distance} & Minimum \termSRD distance used for impulse-like signals. \\
\hline
\code{impulse_like} & \code{epsilon_SRD} & Epsilon term used in the \termSRD denominator for impulse-like signals. \\
\hline
\code{RMS_thresholds} & \code{<signal_name>} & Per-signal RMS threshold used to decide whether two trajectories are sufficiently different. \\
\hline
\end{longtable}
\endgroup

\phantomsection\label{sec:environment-noise-adder}
\subsection*{\texorpdfstring{\titlecode{noise_adder.py}}{noise\_adder.py}}

\code{noise_adder.py} implements the \termSimulator-specific observation-noise profiles.
It must expose:

\begin{itemize}
\item \agentProposedEdit{agent-remove-5ce36589-28c0-5edb-aeef-5c3c6af3816c}{current NOISE_DICT contains explicit low and high profile objects}{\code{NOISE_DICT}: \termSimulator-specific base observation-noise magnitudes for the low-noise regime \code{low}, typically keyed by observable signal.}{}
\agentProposedEdit{agent-remove-1ad4fe81-0409-566e-86d4-49af79132bd2}{NOISE_LEVEL_MULTIPLIER does not exist in the current interface}{The high-noise regime is obtained by multiplying these magnitudes by the shared \code{NOISE_LEVEL_MULTIPLIER["high"]}; \code{none} uses multiplier \code{0}.}{}
\phantomsection\label{def:noise-adder-interface}
\end{itemize}

\begin{DocCode}
def quantify_noise(clean: pd.DataFrame, noisy: pd.DataFrame, reference: pd.DataFrame) -> dict
\end{DocCode}

\agentProposedEdit{agent-remove-58b7e3b1-88c3-5bfb-8a1f-e34062a45273}{the terminology page should contain paper terminology rather than utility pseudocode}{The \code{add_noise(...)} interface is defined in \hyperref[def:add-noise-utility]{Terminology and Common Utilities}.}{}
\agentProposedEdit
  {agent-add-fdb79479-2ff2-5c0c-9906-e0bbc466c0d5}
  {Add this agent-authored block for review.}
  {}
  {The observation-noise implementation contract is documented locally in this environment-profile section through \code{noise_adder.py}, \code{NOISE_DICT}, and the required callable interface.}

\phantomsection\label{sec:environment-detectability-specific}
\subsection*{Optional Implementation Hooks}

\code{detectability_specific_environment.py} exposes:

\begin{DocCode}
def is_detectable(intervention: pd.DataFrame, baseline: pd.DataFrame, run_parameters: dict, intervention_time: float, min_first_diff: float | None) -> str
\end{DocCode}

This hook is an environment-specific veto used by Stage: Cheap Filter Metrics.
If it returns \code{"no"}, the child is marked as not detectable even if the generic RMS checks pass.


\section*{Validation Rules}
\phantomsection\label{sec:environment-validation-rules}

\subsection*{Validation overview}

Validation exists to catch user-authored environment mistakes before they reach expensive MATLAB simulation batches.
The checks are enforced across the metadata build in \code{workflows/simulate/build_metadata.py}, run-graph compilation in \code{workflows/generate_runs/compile_run_graph.py}, simulation preflight in \code{workflows/simulate/run_pending_sims.py}, and cheap-filter validation in \code{workflows/metrics/compute_metrics.py}.

The environment-level checklist includes:

\begin{itemize}
\item \code{description_levels.json::textual_description::<desc_level>} for \termDescLevel in \{\code{high}\} must start with ``by simulating'' when it is appended to the shared prompt text.
\item \code{description_levels.json::internal_naming_to_agent_facing_signal} must be a non-empty object.
\item The keys of \code{description_levels.json::internal_naming_to_agent_facing_signal}, excluding \code{time}, must match \code{experiment_config.json::observable_signals::observable_signals} exactly and in the same order.
\item The signals listed in \code{experiment_config.json::observable_signals::observable_signals} must appear in either \code{generated/metadata.json::simulink_signals_available} or \code{generated/metadata.json::simscape_signals_available}.
\end{itemize}

The full environment-by-environment catalog is collected in the paper appendix, under \hyperref[docstart:appendix_environment]{Appendix: Environment Catalog}.

\end{document}
